% Part: incompleteness
% Chapter: arithmetization-syntax
% Section: coding-formulas

\documentclass[../../../include/open-logic-section]{subfiles}

\begin{document}

\olfileid{inc}{art}{frm}
\olsection{\printtoken{P}{formula} సంకేతీకరణ}

$\fn{Term}(x)$ సంబంధాన్ని ఆదిమ పునరావృత్తంగా నిర్వచించిన తరువాత, దాన్ని
ఉపయోగించి \tetoken{సూత్రాల}{!!{formula}s} కోసం అనురూపమైన $\fn{Frm}(x)$
సంబంధాన్ని ఆదిమ పునరావృత్తంగా నిర్వచించవచ్చు.

\begin{prop}
$x$ ఒక పరమాణు \tetoken{సూత్రానికి}{!!{formula}} గ్యోడెల్ సంఖ్య అయితేనే
నిలిచే $\fn{Atom}(x)$ సంబంధం ఆదిమ పునరావృత్తం.
\end{prop}

\begin{proof}
క్రింది వాటిలో ఒకటి నిలిచినప్పుడే $x$ ఒక పరమాణు
\tetoken{సూత్రానికి}{!!{formula}} గ్యోడెల్ సంఖ్య:
\begin{enumerate}
\item $n$, $j<x$, $z<x$ ఉన్నాయి; ప్రతి $i<n$కి $\fn{Term}((z)_i)$, ఇంకా
  $x=$
\[
\Gn{\Obj P^n_j(} \concat \fn{flatten}(z) \concat \Gn{)}.
\]
\item $z_1,z_2<x$ ఉన్నాయి; $\fn{Term}(z_1)$, $\fn{Term}(z_2)$, ఇంకా
  $x=$
\[
\Gn{{\eq}(} \concat z_1 \concat \Gn{,} \concat z_2 \concat{\Gn{)}}.
\]
  \tagitem{prvFalse}{$x = \Gn{\lfalse}$.}{}
  \tagitem{prvTrue}{$x = \Gn{\ltrue}$.}{}
\end{enumerate}
\end{proof}

\begin{prop}
\ollabel{prop:frm-primrec}
$x$ ఒక \tetoken{సూత్రానికి}{!!a{formula}} గ్యోడెల్ సంఖ్య అయితేనే నిలిచే
$\fn{Frm}(x)$ సంబంధం ఆదిమ పునరావృత్తం.
\end{prop}

\begin{proof}
సంకేతాల క్రమం~$s$ ఒక \tetoken{సూత్రం}{!!a{formula}} కావడానికి అవసరమైనదీ,
సరిపడేదీ ఏమిటంటే, పరమాణు \tetoken{సూత్రాల}{!!{formula}s} నుంచి నిర్మాణ
నియమాల ప్రకారం $s$ ఎలా ఏర్పడిందో నమోదు చేసే $s_0$, \dots,
$s_{k-1}=s$ అనే \tetoken{సూత్రాల}{!!{formula}} నిర్మాణ క్రమం ఉండటం.
ఈ సూత్రం~$s$కు సంకేతసంఖ్య~$x$ అనుకుందాం. ప్రతి $s_i$ యొక్క సంకేతసంఖ్య
$\leq x$; అంతేకాక $k\leq\len{x}$. నిర్మాణ క్రమం
$\tuple{s_0,\dots,s_{k-1}}$.
అందువల్ల, పదాల సందర్భంలోని వాదనలాగే, అటువంటి నిర్మాణ క్రమపు సంకేతసంఖ్యను
$p_{k-1}^{k(x+1)}$తో పరిమితం చేసి, ఈ సంబంధాన్ని పరిమిత పరిమాణీకరణతో
ఆదిమ పునరావృత్తంగా నిర్వచించవచ్చు.
\sourcecorrection{OLTEART-003}{ఆంగ్ల మూలం నిర్మాణ క్రమం మొత్తపు సంకేతసంఖ్య కూడా $s$ యొక్క సంకేతసంఖ్య~$x$కంటే తక్కువని చెబుతుంది; ఇక్కడ ఉపయోగించిన ప్రధాన-సంఖ్యా ఘాత క్రమ సంకేతీకరణలో అది సాధారణంగా నిజం కాదు. పదాల నిరూపణలో ఇచ్చిన సరైన పరిమితి వాదనను ఇక్కడ అనుసరించాం.}
\end{proof}

\begin{prob}
\olref[inc][art][frm]{prop:frm-primrec}కు,
\olref[inc][art][trm]{prop:term-primrec} మొదటి నిరూపణ పద్ధతిలోనే సవివర
నిరూపణ ఇవ్వండి.
\end{prob}

\begin{prop}
\ollabel{prop:freeocc-primrec}
గ్యోడెల్ సంఖ్య~$x$ గల సూత్రంలోని $i$-వ సంకేతం, గ్యోడెల్ సంఖ్య~$z$ గల
చరానికి స్వేచ్ఛా ఘటన అయితేనే నిలిచే
$\fn{FreeOcc}(x,z,i)$ సంబంధం ఆదిమ పునరావృత్తం.
\end{prop}

\begin{proof}
అభ్యాసం.
\end{proof}

\begin{prob}
\olref[inc][art][frm]{prop:freeocc-primrec}ను నిరూపించండి. ఒక
\tetoken{సూత్రం}{!!a{formula}}లోని ఏ ఉపతంత్రి కూడా
\tetoken{ఒక సూత్రమైతే}{!!a{formula}}, అది ఆ సూత్రానికి
ఉప-\tetoken{సూత్రం}{!!{formula}} అవుతుందనే
వాస్తవాన్ని ఉపయోగించవచ్చు.
\end{prob}

\begin{prop}
$x$ ఒక \tetoken{వాక్యానికి}{!!a{sentence}} గ్యోడెల్ సంఖ్య అయితేనే నిలిచే
$\fn{Sent}(x)$ ధర్మం ఆదిమ పునరావృత్తం.
\end{prop}

\begin{proof}
\tetoken{వాక్యం}{!!{sentence}} అంటే \tetoken{చరాల}{!!{variable}s}
స్వేచ్ఛా ఘటనలు లేని \tetoken{సూత్రం}{!!a{formula}}. కాబట్టి
$\fn{Sent}(x)$ నిలవడానికి అవసరమైనదీ, సరిపడేదీ
\begin{multline*}
\bforall{i<\len{x}}{\bforall{z<x}{}}\\
(\bexists{j<z}{z=\Gn{\Obj v_j}} \lif \lnot\fn{FreeOcc}(x,z,i)).
\end{multline*}
\end{proof}


\end{document}
